\section{清初可跳转事件节点}
以下锚点以清初账本的主事件为单位，连接征服、海疆、地方行政与边地治理。每个节点保留核心取证与限度，不将战报、诏令或后出地方叙事混作同一种事实。
\subsection{1640—1649}
\eventanchor{EQ-1644-SHANHAIGUAN}{清军与吴三桂在山海关击败李自成军…}顺治元年四月（1644年5月，换算待核），清军与吴三桂在山海关击败李自成军，打开进入华北的军事通道。核心取证：《清世祖实录》顺治元年四月山海关军报；《清史稿·世祖本纪》；《清初内国史院满文档案》顺治元年四月关宁军务档。编年证词定位国家叙述中的时间，却需同地方或对方材料校验范围。来源保留：原有置信注记：核心战事可靠；吴三桂求援文书原貌、各军兵数和战场过程受清廷与后出叙事塑形
\eventanchor{EQ-1644-BEIJING-ENTRY}{多尔衮率军进入北京并接收明朝京城…}顺治元年五月（1644年6月前后），多尔衮率军进入北京并接收明朝京城的行政与礼制空间。核心取证：《清世祖实录》顺治元年五月；《清史稿·世祖本纪》；《甲申传信录》及北京地方记事。编年证词定位国家叙述中的时间，却需同地方或对方材料校验范围。来源保留：原有置信注记：入城与接收可靠；城内秩序、迎降者范围和掠夺状况须以不同见证互校
\eventanchor{EQ-1644-SHUNZHI-BEIJING}{顺治帝自盛京抵北京，清廷举行定都…}顺治元年九月（1644年10月前后），顺治帝自盛京抵北京，清廷举行定都与皇帝入城仪式。核心取证：《清世祖实录》顺治元年九月；《清史稿·世祖本纪》；《清初内国史院满文档案》北京迁都礼仪档。编年证词定位国家叙述中的时间，却需同地方或对方材料校验范围。来源保留：原有置信注记：日期与仪式主干可靠；官修礼仪记录不能等同各地承认
\eventanchor{EQ-1644-REGENT-ORDERS}{多尔衮以摄政王名义整顿京师官署、…}顺治元年冬（1644年），多尔衮以摄政王名义整顿京师官署、招降明官并发布安民与征发命令。核心取证：《清世祖实录》顺治元年十月至十二月；《清史稿·职官志》；中国第一历史档案馆藏内阁大库顺治元年题本〔京师官署交接〕。编年证词定位国家叙述中的时间，却需同地方或对方材料校验范围。来源保留：原有置信注记：命令发布可靠；接受任官和地方执行不能从诏令直接推出
\eventanchor{EQ-1645-NANJING-TAKEN}{多铎军攻取南京，南明弘光朝的都城…}顺治二年四月（1645年5月前后），多铎军攻取南京，南明弘光朝的都城失守。核心取证：《清世祖实录》顺治二年四月；《清史稿·世祖本纪》；《南明史料》所收弘光朝奏疏与江宁地方记事。编年证词定位国家叙述中的时间，却需同地方或对方材料校验范围。来源保留：原有置信注记：城陷日期可靠；守军数量、降附动机和战后暴力须并读南明与地方材料
\eventanchor{EQ-1645-HONGGUANG-CAPTURE}{弘光帝朱由崧在芜湖一带被俘并被押…}顺治二年五月（1645年），弘光帝朱由崧在芜湖一带被俘并被押往北京。核心取证：《清世祖实录》顺治二年五月；《清史稿·弘光帝本纪》；《明季南略》及江南地方志相关条目。编年证词定位国家叙述中的时间，却需同地方或对方材料校验范围。来源保留：原有置信注记：被俘及押解主干可靠；参与者责任和移交地点在纪传、地方记事间有异
\eventanchor{EQ-1645-HAIR-ORDER}{清廷重申薙发易服命令，在江南等地…}顺治二年六月（1645年），清廷重申薙发易服命令，在江南等地引发服从、逃亡与抵抗。核心取证：《清世祖实录》顺治二年六月；《清史稿·舆服志》；内阁大库顺治二年兵科题本与江南府县志薙发条。编年证词定位国家叙述中的时间，却需同地方或对方材料校验范围。来源保留：原有置信注记：命令文本可靠；各地执行速度和社会反应高度不一，不能以诏书概括全国
\eventanchor{EQ-1645-YANGZHOU}{扬州攻陷后发生大规模杀戮与掳掠，…}顺治二年四—五月（1645年），扬州攻陷后发生大规模杀戮与掳掠，规模和持续时间存在重大争议。核心取证：《清世祖实录》顺治二年四月扬州军报；《清史稿·史可法传》；王秀楚《扬州十日记》与《扬州府志》。编年证词定位国家叙述中的时间，却需同地方或对方材料校验范围。来源保留：原有置信注记：扬州失陷可靠；死亡数字和“十日”叙事须标为后出见证并与地方材料互证
\eventanchor{EQ-1645-JIANGYIN}{江阴在薙发命令和清军围攻中失守，…}顺治二年闰六月至八月（1645年），江阴在薙发命令和清军围攻中失守，地方记忆长期以守城与屠城叙述该事件。核心取证：《清世祖实录》顺治二年江南军报；《清史稿·世祖本纪》；《江阴县志》及《江阴城守纪》。编年证词定位国家叙述中的时间，却需同地方或对方材料校验范围。来源保留：原有置信注记：城陷主干可靠；围城天数、伤亡和人物言行多依赖地方后出文本
\eventanchor{EQ-1645-ZHEJIANG-CONQUEST}{清军进取杭州、绍兴等地，南明鲁王…}顺治二年六月至十月（1645年），清军进取杭州、绍兴等地，南明鲁王政权转入浙东及海上活动。核心取证：《清世祖实录》顺治二年六月至十月；《清史稿·鲁王传》；《浙江通志》与鲁王政权文书辑录。编年证词定位国家叙述中的时间，却需同地方或对方材料校验范围。来源保留：原有置信注记：城市失守与鲁王转移可靠；各府实际控制和海岛联系须逐地核查
\eventanchor{EQ-1646-LONGWU-CAPTURE}{清军在闽北俘杀隆武帝朱聿键，隆武…}顺治三年八月（1646年），清军在闽北俘杀隆武帝朱聿键，隆武朝终结。核心取证：《清世祖实录》顺治三年八月；《清史稿·隆武帝本纪》；《福建通志》与黄道周文集相关记载。编年证词定位国家叙述中的时间，却需同地方或对方材料校验范围。来源保留：原有置信注记：被俘死亡主干可靠；行止路线与地方助力记载不一
\eventanchor{EQ-1646-SHAOWU-GUANGZHOU}{清军攻入广州，绍武帝绍武帝被杀…}顺治三年十一月（1646年），清军攻入广州，绍武帝绍武帝被杀，广州短暂的绍武朝结束。核心取证：《清世祖实录》顺治三年十一月；《清史稿·绍武帝本纪》；《广东通志》及南明文集。编年证词定位国家叙述中的时间，却需同地方或对方材料校验范围。来源保留：原有置信注记：广州失守可靠；绍武与永历两朝冲突细节主要来自后出南明叙事
\eventanchor{EQ-1646-YONGLI-ZHAOQING}{朱由榔在肇庆即位为永历帝，建立延…}永历元年／顺治三年（1646年冬），朱由榔在肇庆即位为永历帝，建立延续至西南的南明朝廷。核心取证：《清世祖实录》顺治三年两广军报；《清史稿·永历帝本纪》；《永历实录》与《肇庆府志》。编年证词定位国家叙述中的时间，却需同地方或对方材料校验范围。来源保留：原有置信注记：即位时间大体可靠；改元与礼仪细节应按永历文书和清廷军报分别处理
\eventanchor{EQ-1648-JIN-SHENGHUAN}{南昌守将金声桓反清并改奉南明，江…}顺治五年（1648年），南昌守将金声桓反清并改奉南明，江西战局重新打开。核心取证：《清世祖实录》顺治五年江西军报；《清史稿·金声桓传》；《南昌府志》与南明永历朝文书。编年证词定位国家叙述中的时间，却需同地方或对方材料校验范围。来源保留：原有置信注记：反叛及南昌控制可靠；部众规模和改奉动机多为交战各方归责语言
\eventanchor{EQ-1648-LI-CHENGDONG}{广东清将李成栋反清，广州一度转入…}顺治五年（1648年），广东清将李成栋反清，广州一度转入永历朝控制。核心取证：《清世祖实录》顺治五年广东军报；《清史稿·李成栋传》；《广东通志》及永历朝诏敕辑录。编年证词定位国家叙述中的时间，却需同地方或对方材料校验范围。来源保留：原有置信注记：反叛与广州易手可靠；个人动机和地方支持范围有强烈阵营叙事差异
\eventanchor{EQ-1649-NANCHANG}{清军围攻并攻取南昌，金声桓反叛失…}顺治六年正月至三月（1649年），清军围攻并攻取南昌，金声桓反叛失败。核心取证：《清世祖实录》顺治六年正月至三月；《清史稿·金声桓传》；内阁大库顺治六年江西军需题本与《南昌府志》。编年证词定位国家叙述中的时间，却需同地方或对方材料校验范围。来源保留：原有置信注记：围城结局可靠；攻城损失、城内处置和战后人口变化须以地方资料复核
\subsection{1650—1659}
\eventanchor{EQ-1650-GUANGZHOU}{清军重新攻取广州，战后杀戮与损毁…}顺治七年十一月（1650年），清军重新攻取广州，战后杀戮与损毁规模成为广东地方记忆中的重大争议。核心取证：《清世祖实录》顺治七年十一月；《清史稿·世祖本纪》；《广东通志》、广州地方笔记与内阁大库广东军报。编年证词定位国家叙述中的时间，却需同地方或对方材料校验范围。来源保留：原有置信注记：广州再陷可靠；死亡数字和“庚寅之劫”持续时间不能从清军战报单定
\eventanchor{EQ-1650-DORGON-DEATH}{摄政王多尔衮在古北口外行猎途中去…}顺治七年十二月（1650年），摄政王多尔衮在古北口外行猎途中去世，摄政政治结构随之改组。核心取证：《清世祖实录》顺治七年十二月；《清史稿·多尔衮传》；《清初内国史院满文档案》摄政王丧仪与遗命档。编年证词定位国家叙述中的时间，却需同地方或对方材料校验范围。来源保留：原有置信注记：死亡日期可靠；死因、身后清算与权力斗争不能只依后来的政治定性
\eventanchor{EQ-1651-SHUNZHI-PERSONAL-RULE}{顺治帝亲政，撤除摄政安排并重置皇…}顺治八年（1651年），顺治帝亲政，撤除摄政安排并重置皇帝周边的决策网络。核心取证：《清世祖实录》顺治八年正月；《清史稿·世祖本纪》；《顺治朝内阁档》亲政诏令与任官题本。编年证词定位国家叙述中的时间，却需同地方或对方材料校验范围。来源保留：原有置信注记：亲政事实可靠；实际决策仍受宗室、旗务和满汉官僚共同制约
\eventanchor{EQ-1652-DINGGUO-GUILIN}{李定国攻克桂林，孔有德自尽，清在…}永历六年／顺治九年（1652年），李定国攻克桂林，孔有德自尽，清在广西的军事据点受重创。核心取证：《清世祖实录》顺治九年广西军报；《清史稿·孔有德传》；《永历实录》与《广西通志》。编年证词定位国家叙述中的时间，却需同地方或对方材料校验范围。来源保留：原有置信注记：桂林战果可靠；双方兵数和孔有德死因细节应保留不同文本说法
\eventanchor{EQ-1652-DINGGUO-HENGZHOU}{李定国在衡州击败尼堪军，清军西南…}永历六年／顺治九年夏（1652年），李定国在衡州击败尼堪军，清军西南攻势暂受挫。核心取证：《清世祖实录》顺治九年湖南军报；《清史稿·尼堪传》；《永历实录》与《衡州府志》。编年证词定位国家叙述中的时间，却需同地方或对方材料校验范围。来源保留：原有置信注记：战役结局大体可靠；地点、伤亡与尼堪死亡过程在中清叙事与南明记录间有异
\eventanchor{EQ-1653-ZHENG-XIAMEN}{郑成功以厦门、金门为基地整合海上…}顺治十年（1653年），郑成功以厦门、金门为基地整合海上军政与贸易网络，持续攻击福建沿海清军据点。核心取证：《清世祖实录》顺治十年福建海防军报；《清史稿·郑成功传》；《厦门志》与郑氏文书辑录。编年证词定位国家叙述中的时间，却需同地方或对方材料校验范围。来源保留：原有置信注记：基地经营可靠；郑军船只、兵员和贸易收入多来自敌对战报或后出郑氏叙事
\eventanchor{EQ-1654-ZHENG-ZHANGZHOU}{郑成功军围攻漳州等闽南城镇，清军…}顺治十一年（1654年），郑成功军围攻漳州等闽南城镇，清军守城与援军维持陆海据点。核心取证：《清世祖实录》顺治十一年福建军报；《清史稿·郑成功传》；《漳州府志》与荷兰东印度公司台湾商馆日记。编年证词定位国家叙述中的时间，却需同地方或对方材料校验范围。来源保留：原有置信注记：围攻与未克主干可靠；郑氏作战目标和双方损失须以中荷航海记录交叉
\eventanchor{EQ-1656-ZHENG-QUANZHOU}{郑成功军在泉州、同安一线进攻并争…}顺治十三年（1656年），郑成功军在泉州、同安一线进攻并争夺闽南水陆通道。核心取证：《清世祖实录》顺治十三年福建军报；《清史稿·郑成功传》；《泉州府志》及郑氏文书辑录。编年证词定位国家叙述中的时间，却需同地方或对方材料校验范围。来源保留：原有置信注记：战事发生可靠；具体攻守日期与城寨归属须按清军报、郑氏文书和地方志逐项核对
\eventanchor{EQ-1657-YONGLI-GUIZHOU}{永历朝在贵州、云南的孙可望与李定…}永历十一年／顺治十四年（1657年），永历朝在贵州、云南的孙可望与李定国集团矛盾公开化，孙可望转降清廷。核心取证：《清世祖实录》顺治十四年贵州军报；《清史稿·孙可望传》；《永历实录》与贵州地方志。编年证词定位国家叙述中的时间，却需同地方或对方材料校验范围。来源保留：原有置信注记：孙可望降清可靠；冲突责任、军队规模与永历帝处境受各方政治文本影响
\eventanchor{EQ-1658-QING-YUNNAN}{清军由贵州、四川等方向进逼云南，…}顺治十五年（1658年），清军由贵州、四川等方向进逼云南，永历朝廷撤离昆明并向西南退却。核心取证：《清世祖实录》顺治十五年云南军报；《清史稿·洪承畴传》；《云南通志》与永历朝奏疏辑录。编年证词定位国家叙述中的时间，却需同地方或对方材料校验范围。来源保留：原有置信注记：清军入滇和永历撤离可靠；各路行军与地方归附应分区而非视作全省一次完成
\eventanchor{EQ-1659-NANJING-EXPEDITION}{郑成功北上围攻南京，未能突破城防…}顺治十六年七月至八月（1659年），郑成功北上围攻南京，未能突破城防后撤回海上基地。核心取证：《清世祖实录》顺治十六年七月至八月；《清史稿·郑成功传》；《延平王户官杨英从征实录》与荷兰东印度公司海上报告。编年证词定位国家叙述中的时间，却需同地方或对方材料校验范围。来源保留：原有置信注记：北伐及撤退可靠；作战计划、内应和伤亡数字在清廷、郑氏与荷兰记录间冲突
\eventanchor{EQ-1659-YONGLI-MYANMAR}{永历帝一行进入缅甸，西南战争扩展…}永历十三年／顺治十六年末（1659年），永历帝一行进入缅甸，西南战争扩展为跨境避难和外交问题。核心取证：《清世祖实录》顺治十六年云南军报；《清史稿·永历帝本纪》；《琉璃宫史》缅文传统及欧洲旅行记译注。编年证词定位国家叙述中的时间，却需同地方或对方材料校验范围。来源保留：原有置信注记：入缅事实可靠；路线、接待条件与缅甸朝廷决策须并读汉文和缅文/欧洲记载
\subsection{1660—1669}
\eventanchor{EQ-1660-QING-YUNNAN-CONSOLIDATION}{清军继续经营云南城镇与驿路，永历…}顺治十七年（1660年），清军继续经营云南城镇与驿路，永历朝残部向中缅边地退缩。核心取证：《清世祖实录》顺治十七年云南军报；《清史稿·吴三桂传》；内阁大库顺治十七年云南粮饷题本与《云南通志》。编年证词定位国家叙述中的时间，却需同地方或对方材料校验范围。来源保留：原有置信注记：清军控制主要城镇大体可靠；土司、矿区和山地实际控制不能以府城易手推定
\eventanchor{EQ-1661-KANGXI-ACCESSION}{顺治帝去世，八岁玄烨即位为康熙帝…}顺治十八年正月（1661年），顺治帝去世，八岁玄烨即位为康熙帝，索尼、苏克萨哈、遏必隆、鳌拜辅政。核心取证：《清世祖实录》顺治十八年正月；《清史稿·圣祖本纪》；《康熙起居注》与内国史院遗诏档。编年证词定位国家叙述中的时间，却需同地方或对方材料校验范围。来源保留：原有置信注记：继位及辅政名单可靠；顺治死因和遗诏形成过程存在材料与后世叙事差异
\eventanchor{EQ-1661-COAST-EVACUATION}{清廷下令福建广东沿海迁界，试图切…}顺治十八年（1661年），清廷下令福建广东沿海迁界，试图切断郑氏的粮食、人口和贸易联系。核心取证：《清世祖实录》顺治十八年迁界谕旨；《清史稿·食货志》；《福建通志》《广东通志》及沿海县志迁界条。编年证词定位国家叙述中的时间，却需同地方或对方材料校验范围。来源保留：原有置信注记：迁界命令可靠；实际迁撤距离、死亡与秘密贸易须以各县档案和方志核实
\eventanchor{EQ-1661-ZHENG-TAIWAN-LANDING}{郑成功军登陆台湾并围攻荷兰东印度…}永历十五年／顺治十八年四月（1661年），郑成功军登陆台湾并围攻荷兰东印度公司在热兰遮城的据点。核心取证：《清世祖实录》顺治十八年海防军报；《清史稿·郑成功传》；《热兰遮城日志》荷文1661年四月至十二月条。编年证词定位国家叙述中的时间，却需同地方或对方材料校验范围。来源保留：原有置信注记：登陆与围城可靠；兵员、原住民联盟与双方伤亡需区分中荷各自战报
\eventanchor{EQ-1662-ZEELANDIA}{热兰遮城守军投降，荷兰东印度公司…}永历十六年／康熙元年正月（1662年），热兰遮城守军投降，荷兰东印度公司在台湾西南的统治终结。核心取证：《清圣祖实录》康熙元年台湾相关军报；《清史稿·郑成功传》；《热兰遮城日志》荷文1662年正月及投降条约。编年证词定位国家叙述中的时间，却需同地方或对方材料校验范围。来源保留：原有置信注记：投降文本和日期可靠；城外居民、原住民与被俘人员经历须从不同语言材料重建
\eventanchor{EQ-1662-YONGLI-EXECUTION}{吴三桂军将自缅甸押回的永历帝朱由…}康熙元年四月（1662年），吴三桂军将自缅甸押回的永历帝朱由榔处死，南明朝廷名义中心结束。核心取证：《清圣祖实录》康熙元年四月；《清史稿·永历帝本纪》；《琉璃宫史》缅文传统与《云南通志》。编年证词定位国家叙述中的时间，却需同地方或对方材料校验范围。来源保留：原有置信注记：被押回和处死主干可靠；缅甸移交过程、处死地点与责任叙事须比对中缅材料
\eventanchor{EQ-1662-ZHENG-DEATH}{郑成功在台湾去世，郑氏集团进入继…}永历十六年／康熙元年五月（1662年），郑成功在台湾去世，郑氏集团进入继承与军政重组阶段。核心取证：《清圣祖实录》康熙元年海防军报；《清史稿·郑成功传》；《巴达维亚城日记》1662年台湾消息与郑氏文书辑录。编年证词定位国家叙述中的时间，却需同地方或对方材料校验范围。来源保留：原有置信注记：死亡月份大体可靠；死因、继承人之争与军中处置依赖郑氏后出材料和荷兰报告
\eventanchor{EQ-1663-XIAMEN-JINMEN}{清军联合部分郑氏降将攻取厦门、金…}康熙二年（1663年），清军联合部分郑氏降将攻取厦门、金门，郑氏主力退守台湾。核心取证：《清圣祖实录》康熙二年福建军报；《清史稿·施琅传》；《巴达维亚城日记》1663年中国沿海消息。编年证词定位国家叙述中的时间，却需同地方或对方材料校验范围。来源保留：原有置信注记：厦金易手可靠；荷兰参与、海战细节和郑氏损失应区分清方军报与荷兰记录
\eventanchor{EQ-1664-KUIZHOU}{清军攻破夔东据点，李来亨等抗清武…}康熙三年（1664年），清军攻破夔东据点，李来亨等抗清武装失败，长江上游的明遗民军事中心瓦解。核心取证：《清圣祖实录》康熙三年湖广四川军报；《清史稿·李来亨传》；《夔州府志》与南明遗民文集。编年证词定位国家叙述中的时间，却需同地方或对方材料校验范围。来源保留：原有置信注记：夔东结局可靠；各部关系、山地据点与居民伤亡材料零散
\eventanchor{EQ-1665-REGENTS-CONSOLIDATION}{四辅政大臣以议政和旗务网络处理军…}康熙四年（1665年），四辅政大臣以议政和旗务网络处理军政事务，鳌拜在政务中的影响扩大。核心取证：《清圣祖实录》康熙四年政务条；《清史稿·鳌拜传》；《康熙起居注》及内阁大库满汉题本。编年证词定位国家叙述中的时间，却需同地方或对方材料校验范围。来源保留：原有置信注记：辅政架构可靠；权力高下多由后来的鳌拜案档与《清史稿》倒叙，须避免预设结局
\eventanchor{EQ-1667-SONI-DEATH}{首辅索尼去世，四辅政大臣的均衡被…}康熙六年（1667年），首辅索尼去世，四辅政大臣的均衡被打破。核心取证：《清圣祖实录》康熙六年；《清史稿·索尼传》；《康熙起居注》康熙六年丧仪与议政记录。编年证词定位国家叙述中的时间，却需同地方或对方材料校验范围。来源保留：原有置信注记：索尼死亡及官职变动可靠；其政治立场与遗言多经后出人物传记修饰
\eventanchor{EQ-1669-OBOI-ARREST}{康熙帝拘捕鳌拜并审理其党羽，结束…}康熙八年五月（1669年），康熙帝拘捕鳌拜并审理其党羽，结束辅政大臣实际主政的局面。核心取证：《清圣祖实录》康熙八年五月；《清史稿·鳌拜传》；《康熙起居注》与鳌拜案内阁档。编年证词定位国家叙述中的时间，却需同地方或对方材料校验范围。来源保留：原有置信注记：拘捕与处分可靠；宫中擒捕过程、罪状与少年皇帝个人谋略含胜利者叙事
\subsection{1670—1679}
\eventanchor{EQ-1670-SACRED-EDICT}{康熙帝颁布十六条圣谕，规定由地方…}康熙九年（1670年），康熙帝颁布十六条圣谕，规定由地方官和民众遵循的伦理、赋役与秩序训导框架。核心取证：《清圣祖实录》康熙九年十月；《清史稿·圣祖本纪》；《大清会典事例》圣谕条与各省地方志宣讲记录。编年证词定位国家叙述中的时间，却需同地方或对方材料校验范围。来源保留：原有置信注记：颁布可靠；宣讲频率、识读范围和实际社会效果不能由文本证明
\eventanchor{EQ-1673-FEUDATORY-ABOLITION}{康熙帝批准撤除吴三桂、尚可喜、耿…}康熙十二年三月（1673年），康熙帝批准撤除吴三桂、尚可喜、耿继茂三藩的请求或安排，触发藩镇利益危机。核心取证：《清圣祖实录》康熙十二年三月；《清史稿·三藩传》；内阁大库康熙十二年撤藩题本与《平定三逆方略》。编年证词定位国家叙述中的时间，却需同地方或对方材料校验范围。来源保留：原有置信注记：撤藩决策可靠；各藩请求先后、朝议立场和财政核算须以奏折及部议分开处理
\eventanchor{EQ-1673-WU-REVOLT}{吴三桂在云南起兵反清，三藩之乱正…}康熙十二年十一月（1673年），吴三桂在云南起兵反清，三藩之乱正式爆发。核心取证：《清圣祖实录》康熙十二年十一月；《清史稿·吴三桂传》；《平定三逆方略》卷首及云南地方志。编年证词定位国家叙述中的时间，却需同地方或对方材料校验范围。来源保留：原有置信注记：起兵及时间可靠；檄文中的族群、忠明和财政诉求是动员语言，不能直接等同参与者动机
\eventanchor{EQ-1674-GENG-JINGZHONG}{靖南王耿精忠在福建反清，东南海陆…}康熙十三年（1674年），靖南王耿精忠在福建反清，东南海陆战场与郑氏关系重新组合。核心取证：《清圣祖实录》康熙十三年福建军报；《清史稿·耿精忠传》；《平定三逆方略》与《福建通志》。编年证词定位国家叙述中的时间，却需同地方或对方材料校验范围。来源保留：原有置信注记：耿精忠反叛可靠；其与郑氏盟约和实际协同行动须逐项辨析
\eventanchor{EQ-1674-WANG-FUCHEN}{陕西提督王辅臣据平凉反清，西北军…}康熙十三年（1674年），陕西提督王辅臣据平凉反清，西北军政和通往甘肃的交通受威胁。核心取证：《清圣祖实录》康熙十三年陕西军报；《清史稿·王辅臣传》；《平定三逆方略》与《平凉府志》。编年证词定位国家叙述中的时间，却需同地方或对方材料校验范围。来源保留：原有置信注记：平凉反叛可靠；王辅臣与吴三桂联络程度、兵数和民众态度不宜由檄文推定
\eventanchor{EQ-1674-ZHENG-JING-FUJIAN}{郑经自台湾出兵福建，取得漳泉沿海…}康熙十三年（1674年），郑经自台湾出兵福建，取得漳泉沿海若干据点并介入三藩战争。核心取证：《清圣祖实录》康熙十三年福建海防军报；《清史稿·郑成功传附郑经》；《台湾外纪》与《福建通志》。编年证词定位国家叙述中的时间，却需同地方或对方材料校验范围。来源保留：原有置信注记：出兵和若干据点易手可靠；郑耿关系、控制范围和征粮情况应以地方志与荷兰/商船信息校核
\eventanchor{EQ-1675-SHANG-ZHIXIN}{平南王尚之信在广州控制军政并与反…}康熙十四年（1675年），平南王尚之信在广州控制军政并与反清势力周旋，两广局势动荡。核心取证：《清圣祖实录》康熙十四年广东军报；《清史稿·尚之信传》；《平定三逆方略》与《广东通志》。编年证词定位国家叙述中的时间，却需同地方或对方材料校验范围。来源保留：原有置信注记：尚之信的反复与广州控制大体可靠；是否完全反清须按不同时段命令、行动和归降分别记录
\eventanchor{EQ-1676-GENG-SURRENDER}{耿精忠向清廷投降，福建藩镇叛乱的…}康熙十五年（1676年），耿精忠向清廷投降，福建藩镇叛乱的核心力量瓦解。核心取证：《清圣祖实录》康熙十五年；《清史稿·耿精忠传》；内阁大库康熙十五年福建招抚题本与《福建通志》。编年证词定位国家叙述中的时间，却需同地方或对方材料校验范围。来源保留：原有置信注记：投降与处分可靠；部众去向、郑氏影响和福建地方恢复须区别官方赦令与实际情况
\eventanchor{EQ-1676-WUCHANG}{清军收复武昌等长江中游要地，吴军…}康熙十五年（1676年），清军收复武昌等长江中游要地，吴军向湖南、广西方向收缩。核心取证：《清圣祖实录》康熙十五年湖广军报；《清史稿·赵良栋传》；《平定三逆方略》与《武昌府志》。编年证词定位国家叙述中的时间，却需同地方或对方材料校验范围。来源保留：原有置信注记：武昌易手可靠；战役阶段、民众遭遇和双方控制线需按府县而非一城外推
\eventanchor{EQ-1677-WU-ZHOU}{吴三桂在衡阳称帝，国号周，将撤藩…}康熙十六年（1677年），吴三桂在衡阳称帝，国号周，将撤藩战争转为公开争夺最高政治名义。核心取证：《清圣祖实录》康熙十六年湖南军报；《清史稿·吴三桂传》；《平定三逆方略》及吴周檄文辑录。编年证词定位国家叙述中的时间，却需同地方或对方材料校验范围。来源保留：原有置信注记：称帝可靠；礼制、官号与各地响应范围主要见敌对军报，需与吴周文书互证
\eventanchor{EQ-1678-WU-DEATH}{吴三桂去世，吴世璠继承吴周政权}康熙十七年八月（1678年），吴三桂去世，吴世璠继承吴周政权。核心取证：《清圣祖实录》康熙十七年八月；《清史稿·吴三桂传》；《平定三逆方略》与湖南地方志。编年证词定位国家叙述中的时间，却需同地方或对方材料校验范围。来源保留：原有置信注记：死亡与继承可靠；死因、遗命和军中稳定程度多来自后出叙事，应避免戏剧化细节
\eventanchor{EQ-1679-HUNAN-RECOVERY}{清军在湖南持续推进并夺回若干府县…}康熙十八年（1679年），清军在湖南持续推进并夺回若干府县，吴周军向贵州、广西收缩。核心取证：《清圣祖实录》康熙十八年湖南军报；《清史稿·赵良栋传》；《平定三逆方略》及《湖南通志》。编年证词定位国家叙述中的时间，却需同地方或对方材料校验范围。来源保留：原有置信注记：清军推进主干可靠；具体县城归属和居民逃亡须由地方志、契约和军报逐县核查
\subsection{1680—1689}
\eventanchor{EQ-1680-SICHUAN-RECOVERY}{清军清理四川通往云南贵州的要道并…}康熙十九年（1680年），清军清理四川通往云南贵州的要道并恢复主要府城控制。核心取证：《清圣祖实录》康熙十九年四川军报；《清史稿·傅弘烈传》；《四川通志》与内阁大库四川粮饷题本。编年证词定位国家叙述中的时间，却需同地方或对方材料校验范围。来源保留：原有置信注记：主要军事结果可靠；四川战后人口、耕地和流亡统计口径极不一致
\eventanchor{EQ-1681-KUNMING}{清军攻入昆明，吴世璠自尽，三藩之…}康熙二十年十月（1681年），清军攻入昆明，吴世璠自尽，三藩之乱的吴周政权终结。核心取证：《清圣祖实录》康熙二十年十月；《清史稿·吴世璠传》；《平定三逆方略》与《云南通志》。编年证词定位国家叙述中的时间，却需同地方或对方材料校验范围。来源保留：原有置信注记：昆明失守和吴世璠死亡可靠；围城损失、降附与后续清剿不得视为同日完成
\eventanchor{EQ-1681-ZHENG-JING-DEATH}{郑经去世，东宁发生继承冲突，郑克…}康熙二十年（1681年），郑经去世，东宁发生继承冲突，郑克塽最终掌权。核心取证：《清圣祖实录》康熙二十年台湾情报；《清史稿·郑成功传附郑经》；《台湾外纪》与荷兰商馆残存通信。编年证词定位国家叙述中的时间，却需同地方或对方材料校验范围。来源保留：原有置信注记：郑经死期和继承结果可靠；宫廷政变过程主要依郑氏后出材料与清方情报
\eventanchor{EQ-1682-SHI-LANG}{清廷任命施琅统率福建水师，筹备进…}康熙二十一年（1682年），清廷任命施琅统率福建水师，筹备进攻台湾。核心取证：《清圣祖实录》康熙二十一年福建海防条；《清史稿·施琅传》；《康熙朝满文朱批奏折全译》康熙二十一年台湾军务奏折。编年证词定位国家叙述中的时间，却需同地方或对方材料校验范围。来源保留：原有置信注记：任命和筹备可靠；施琅战略构想与朝廷内部争论须以奏折、朱批和后出传记区别处理
\eventanchor{EQ-1683-PENGHU}{施琅水师在澎湖击败郑氏舰队，取得…}康熙二十二年六月（1683年），施琅水师在澎湖击败郑氏舰队，取得通往台湾本岛的制海优势。核心取证：《清圣祖实录》康熙二十二年六月；《清史稿·施琅传》；施琅《飞报澎湖大捷疏》与台湾地方志。编年证词定位国家叙述中的时间，却需同地方或对方材料校验范围。来源保留：原有置信注记：澎湖战役结果可靠；舰船、兵员和死伤数字多来自胜方战报，须以郑氏和海域材料约束
\eventanchor{EQ-1683-TAIWAN-SURRENDER}{郑克塽向清廷投降，东宁政权结束}康熙二十二年八月（1683年），郑克塽向清廷投降，东宁政权结束。核心取证：《清圣祖实录》康熙二十二年八月；《清史稿·郑成功传附郑克塽》；施琅奏折与《台湾府志》。编年证词定位国家叙述中的时间，却需同地方或对方材料校验范围。来源保留：原有置信注记：投降与清廷接收主干可靠；投降条件、官兵安置和原住民、移民社会变化须分层记录
\eventanchor{EQ-1684-TAIWAN-PREFECTURE}{清廷设置台湾府及诸县，隶属福建省…}康熙二十三年（1684年），清廷设置台湾府及诸县，隶属福建省，开始以府县和驻防制度经营台湾。核心取证：《清圣祖实录》康熙二十三年台湾设治条；《清史稿·地理志》；《台湾府志》与福建省档案相关奏议。编年证词定位国家叙述中的时间，却需同地方或对方材料校验范围。来源保留：原有置信注记：设府设县可靠；机构设置不等于全岛有效治理，原住民地区和内山控制尤须保留
\eventanchor{EQ-1684-SEA-BAN-LIFT}{清廷调整迁界与海禁政策，允许沿海…}康熙二十三年（1684年），清廷调整迁界与海禁政策，允许沿海居民复界并逐步恢复民间海贸。核心取证：《清圣祖实录》康熙二十三年复界海禁条；《清史稿·食货志》；《福建通志》《广东通志》沿海复界条。编年证词定位国家叙述中的时间，却需同地方或对方材料校验范围。来源保留：原有置信注记：政策转向可靠；复界日期、港口开放和走私变化因省县而异，不能由中央谕旨推全国
\eventanchor{EQ-1685-ALBAZIN-FIRST}{清军围攻并迫使俄方撤离阿尔巴津，…}康熙二十四年（1685年），清军围攻并迫使俄方撤离阿尔巴津，黑龙江边境进入直接谈判前的军事对峙。核心取证：《清圣祖实录》康熙二十四年黑龙江军报；《清史稿·邦交志》；俄国西伯利亚衙门文件与尼布楚谈判材料。编年证词定位国家叙述中的时间，却需同地方或对方材料校验范围。来源保留：原有置信注记：围城与撤离可靠；守军人数、伤亡和领土主张分别来自清俄两方，不能以一方战报定量
\eventanchor{EQ-1686-ALBAZIN-SECOND}{俄方重返阿尔巴津后，清军再次围困…}康熙二十五年（1686年），俄方重返阿尔巴津后，清军再次围困，双方在疾病、补给与外交压力下僵持。核心取证：《清圣祖实录》康熙二十五年黑龙江军报；《清史稿·邦交志》；俄国阿尔巴津驻军报告及城址考古报告。编年证词定位国家叙述中的时间，却需同地方或对方材料校验范围。来源保留：原有置信注记：再围困可靠；兵力、疾病死亡与停战时点应并读满文、俄文与考古材料
\eventanchor{EQ-1688-GALDAN-KHALKHA}{噶尔丹进攻喀尔喀，喀尔喀诸部南徙…}康熙二十七年（1688年），噶尔丹进攻喀尔喀，喀尔喀诸部南徙并向清廷求援，北方草原秩序改变。核心取证：《清圣祖实录》康熙二十七年蒙古军报；《清史稿·噶尔丹传》；《蒙古源流》与俄罗斯使团材料。编年证词定位国家叙述中的时间，却需同地方或对方材料校验范围。来源保留：原有置信注记：入侵与迁徙主干可靠；部众规模、附属关系和“归附”表述须并读蒙文与满汉文
\eventanchor{EQ-1689-NERCHINSK}{清俄在尼布楚签订条约，以满文、拉…}康熙二十八年七月（1689年），清俄在尼布楚签订条约，以满文、拉丁文和俄文文本处理边界与阿尔巴津问题。核心取证：《清圣祖实录》康熙二十八年七月；《清史稿·邦交志》；《尼布楚条约》满文、拉丁文、俄文文本及俄国谈判档。编年证词定位国家叙述中的时间，却需同地方或对方材料校验范围。来源保留：原有置信注记：签约和多文本存在可靠；各文本术语、地名对应和界线含义不完全一致，不能投射为精确现代边界
\subsection{1690—1699}
\eventanchor{EQ-1690-ULAN-BUTUNG}{清军在乌兰布通与噶尔丹军交战，阻…}康熙二十九年（1690年），清军在乌兰布通与噶尔丹军交战，阻止其向内蒙古腹地深入但未立即消灭其军队。核心取证：《清圣祖实录》康熙二十九年八月；《清史稿·噶尔丹传》；《亲征平定朔漠方略》及俄国商旅报告。编年证词定位国家叙述中的时间，却需同地方或对方材料校验范围。来源保留：原有置信注记：战役与战略结果大体可靠；胜负判断、兵力和伤亡在方略与准噶尔相关材料中不一
\eventanchor{EQ-1691-DOLONNOR}{康熙帝在多伦诺尔会盟，安排喀尔喀…}康熙三十年（1691年），康熙帝在多伦诺尔会盟，安排喀尔喀诸部入盟旗体系并调度对噶尔丹资源。核心取证：《清圣祖实录》康熙三十年五月；《清史稿·藩部传》；理藩院满蒙文档案中的喀尔喀盟旗与赈济文书。编年证词定位国家叙述中的时间，却需同地方或对方材料校验范围。来源保留：原有置信注记：会盟和册封可靠；“归附”不是单向完成，盟旗编制、牧地与领主自主性须以蒙文档案追踪
\eventanchor{EQ-1695-GALDAN-CAMPAIGN}{康熙帝筹划并发动对噶尔丹的第二阶…}康熙三十四年（1695年），康熙帝筹划并发动对噶尔丹的第二阶段远征，集中军粮、驼马和多路部队。核心取证：《清圣祖实录》康熙三十四年军务条；《清史稿·圣祖本纪》；《康熙朝满文朱批奏折全译》噶尔丹军需奏折。编年证词定位国家叙述中的时间，却需同地方或对方材料校验范围。来源保留：原有置信注记：远征部署可靠；实际行军人数、牲畜损耗和地方征发须以军需档而非方略夸示数字判断
\eventanchor{EQ-1696-JAOMODO}{清军在昭莫多一带击败噶尔丹主力，…}康熙三十五年（1696年），清军在昭莫多一带击败噶尔丹主力，噶尔丹向西遁走。核心取证：《清圣祖实录》康熙三十五年五月；《清史稿·噶尔丹传》；《亲征平定朔漠方略》与俄国边境报告。编年证词定位国家叙述中的时间，却需同地方或对方材料校验范围。来源保留：原有置信注记：战役结局可靠；会战位置、伤亡与各部贡献应并读满文军报、蒙古材料和俄方消息
\eventanchor{EQ-1697-GALDAN-DEATH}{噶尔丹在科布多一带去世，其部众分…}康熙三十六年（1697年），噶尔丹在科布多一带去世，其部众分散或向准噶尔核心地区退却。核心取证：《清圣祖实录》康熙三十六年；《清史稿·噶尔丹传》；俄国西伯利亚档案与蒙古口传整理。编年证词定位国家叙述中的时间，却需同地方或对方材料校验范围。来源保留：原有置信注记：死亡事实大体可靠；死因、最后地点和部众归属说法各异，清廷凯旋叙事须与俄蒙材料互校
\eventanchor{EQ-1698-KHALKHA-ADMINISTRATION}{清廷在噶尔丹败亡后调整喀尔喀盟旗…}康熙三十七年（1698年），清廷在噶尔丹败亡后调整喀尔喀盟旗、台站与赈济安排，将战时同盟转为日常治理。核心取证：《清圣祖实录》康熙三十七年蒙古条；《清史稿·藩部传》；理藩院满蒙文档案中的喀尔喀盟旗与赈济文书。编年证词定位国家叙述中的时间，却需同地方或对方材料校验范围。来源保留：原有置信注记：制度调整可由谕令和理藩院文书确认；实际牧地使用、移徙和领主关系需长期追踪
\subsection{1700—1709}
\eventanchor{EQ-1705-LHAZANG-TAKEOVER}{拉藏汗推翻第巴桑结嘉措的政权安排…}康熙四十四年（1705年），拉藏汗推翻第巴桑结嘉措的政权安排，西藏政教权力结构剧烈变化。核心取证：《清圣祖实录》康熙四十四年西藏条；《清史稿·西藏传》；《西藏王臣记》藏文与清廷驻藏来文。编年证词定位国家叙述中的时间，却需同地方或对方材料校验范围。来源保留：原有置信注记：政变主干可靠；清廷知情、支持程度和地方反应须并读藏文、满文与旅行者材料
\eventanchor{EQ-1706-SIXTH-DALAI}{第六世达赖喇嘛仓央嘉措被押离拉萨…}康熙四十五年（1706年），第六世达赖喇嘛仓央嘉措被押离拉萨，途中死亡或失踪，其结局与继承合法性成为争议。核心取证：《清圣祖实录》康熙四十五年西藏条；《清史稿·西藏传》；第巴传记、藏文史书与拉达克、蒙古相关记载。编年证词定位国家叙述中的时间，却需同地方或对方材料校验范围。来源保留：原有置信注记：被押离拉萨可靠；死亡地点、日期和是否在世的传统说法相互冲突，不能写成确定结论
\eventanchor{EQ-1708-KELZANG}{卡尔藏嘉措在青海塔尔寺附近被认定…}康熙四十七年（1708年），卡尔藏嘉措在青海塔尔寺附近被认定为第六世达赖喇嘛转世之一，围绕其身份形成跨区域支持网络。核心取证：《清圣祖实录》康熙四十七年青海西藏条；《清史稿·西藏传》；塔尔寺藏文档案与第七世达赖喇嘛传。编年证词定位国家叙述中的时间，却需同地方或对方材料校验范围。来源保留：原有置信注记：出生与早期认定大体可靠；具体认定程序、时间和各方承认不可由后世转世谱系倒推
\subsection{1710—1719}
\eventanchor{EQ-1711-TAX-FREEZE}{康熙帝宣布以康熙五十年丁数为定额…}康熙五十年（1711年），康熙帝宣布以康熙五十年丁数为定额，后生人丁不再加征丁税。核心取证：《清圣祖实录》康熙五十年丁银谕；《清史稿·食货志》；户部档案与《大清会典事例》赋役条。编年证词定位国家叙述中的时间，却需同地方或对方材料校验范围。来源保留：原有置信注记：诏令可靠；各省执行时间、丁银与地丁并征口径和实际负担不能由诏令推定
\eventanchor{EQ-1712-TAX-EDICT}{清廷再次谕令固定丁额并要求地方遵…}康熙五十一年（1712年），清廷再次谕令固定丁额并要求地方遵行，围绕丁银征收与造册形成持续行政问题。核心取证：《清圣祖实录》康熙五十一年赋役谕；《清史稿·食货志》；户部题本与江苏、直隶等省赋役档。编年证词定位国家叙述中的时间，却需同地方或对方材料校验范围。来源保留：原有置信注记：重复申谕可靠；它证明中央规范而非各地立即停止加派，需追踪地方账册与奏报
\eventanchor{EQ-1717-DZUNGAR-LHASA}{准噶尔军进入拉萨并杀拉藏汗，拉萨…}康熙五十六年（1717年），准噶尔军进入拉萨并杀拉藏汗，拉萨政局再次易手。核心取证：《清圣祖实录》康熙五十六年西藏军报；《清史稿·西藏传》；藏文《西藏王臣记》续编与准噶尔相关材料。编年证词定位国家叙述中的时间，却需同地方或对方材料校验范围。来源保留：原有置信注记：入拉萨和拉藏汗死亡可靠；暴力规模、民众态度与准噶尔治理需并读藏文、满文和旅行记录
\eventanchor{EQ-1718-QING-TIBET-EXPEDITION}{清军自四川方向入藏，在喀喇乌苏一…}康熙五十七年（1718年），清军自四川方向入藏，在喀喇乌苏一带覆败，首批救援行动受挫。核心取证：《清圣祖实录》康熙五十七年川藏军报；《清史稿·西藏传》；《康熙朝满文朱批奏折全译》入藏军需奏折及藏文地方记载。编年证词定位国家叙述中的时间，却需同地方或对方材料校验范围。来源保留：原有置信注记：远征失利可靠；路线、伤亡和藏地向导问题主要见清方军报，需与藏文地名和地方传统核对
\subsection{1720—1729}
\eventanchor{EQ-1720-LHASA-RESTORATION}{清军及青海蒙古盟军进入拉萨，驱逐…}康熙五十九年（1720年），清军及青海蒙古盟军进入拉萨，驱逐准噶尔力量并迎立卡尔藏嘉措为第七世达赖喇嘛。核心取证：《清圣祖实录》康熙五十九年西藏军报；《清史稿·西藏传》；满文军机前身档、藏文第七世达赖传与青海蒙古文书。编年证词定位国家叙述中的时间，却需同地方或对方材料校验范围。来源保留：原有置信注记：进入拉萨和迎立主干可靠；清军、蒙古盟军与西藏地方行动者角色及伤亡不能由凯旋叙事抹平
\eventanchor{EQ-1721-ZHU-YIGUI}{朱一贵在台湾发动反清起事，一度攻…}康熙六十年（1721年），朱一贵在台湾发动反清起事，一度攻入府城，清廷随后派兵恢复控制。核心取证：《清圣祖实录》康熙六十年台湾军报；《清史稿·朱一贵传》；《台湾府志》、福建督抚奏折与契约文书。编年证词定位国家叙述中的时间，却需同地方或对方材料校验范围。来源保留：原有置信注记：起事、府城易手与镇压主干可靠；参与群体、赋税原因和伤亡数字受官府“逆匪”分类影响
\eventanchor{EQ-1722-KANGXI-DEATH}{康熙帝去世，皇四子胤禛即位为雍正…}康熙六十一年十一月（1722年），康熙帝去世，皇四子胤禛即位为雍正帝，早期清朝的长期统治转入新的继承阶段。核心取证：《清圣祖实录》康熙六十一年十一月；《清史稿·圣祖本纪、世宗本纪》；《雍正朝起居注册》与内阁大库遗诏档。编年证词定位国家叙述中的时间，却需同地方或对方材料校验范围。来源保留：原有置信注记：去世与即位日期可靠；遗诏、传位过程和储位争议受政治清算与后世传说影响，须以档案版本互校
